%Dokument pro XeLaTeX, překlad pomocí XeLaTeX
\documentclass[a4paper,oneside,11pt,usenames,dvipsnames]{article}

%XeLaTeXové věci a fonty
\usepackage{fontspec,xltxtra,xunicode}
\defaultfontfeatures{Scale=MatchLowercase}

%Písma
%\setromanfont[Mapping=tex-text]{Linux Libertine O}
%\setsansfont[Mapping=tex-text]{Linux Biolinum O}
%\setmonofont[Mapping=tex-text]{Linux Libertine Mono O}

%Pro správné dělení českých slov
\usepackage{polyglossia}
\setdefaultlanguage{czech} %TODO: Případně změňte jazyk
%České uvozovky buď přímo nebo jako ,,uvozovky``

%Některé speciální znaky
\usepackage{textcomp}
% Url
\usepackage{url}
\def\UrlFont{\ttfamily\slshape}

%Barvy pro hyperref
\usepackage{color}
% Hyperref pro odkazy v PDF a URL
\usepackage{hyperref}
% Hyperref a PDF nastavení
\hypersetup{
    pdftitle={Titulek},    % title TODO: Vyplnit titulek
    pdfauthor={Nějaký Autor},     % author TODO: Vyplnit autora
    colorlinks=true,       % false: boxed links; true: colored links
    linkcolor=RoyalBlue,          % color of internal links
    citecolor=RoyalBlue,        % color of links to bibliography
    filecolor=RoyalBlue,      % color of file links
    urlcolor=RoyalBlue           % color of external links
}

%Nadpisy, záhlaví a zápatí
\usepackage[pagestyles]{titlesec}

%Nadpisy bezpatkovým písmem
\titleformat{\part}[display]{\filright}{\Large\sffamily Část\ \thepart}{0.25em}{\bfseries\sffamily\huge\upshape}
\titleformat{\section}[block]{\sffamily\Large\bfseries}{\thesection}{1em}{\Large}
\titleformat{\subsection}[block]{\sffamily\large\bfseries}{\thesubsection}{1em}{\large}
\titleformat{\subsubsection}[block]{\sffamily\normalsize\bfseries}{\thesubsubsection}{1em}{\normalsize}

\begin{document}

\thispagestyle{empty}
\begin {center}
	\vspace *{10mm}
	{\Huge \sffamily \bfseries Znalostní ontologie z oboru knih}

	\vspace *{20mm}
	{\Large \sffamily \bfseries Ondřej Maxa, Jakub Vojtko, Jakub Vítovec}

	\vspace *{7.5mm}
	{\large \sffamily \today}
\end {center}

\begin{abstract}
    Tato semestrální práce se zabývá analýzou, návrhem, implementací a kritickou analýzou znalostní ontologie zaměřené na obor literatury. Dle metody 101 pro tvorbu ontologií byla vytvořena základní ontologie, zabírající 
    zmíněný obor a následně tato byla implementována v open-source frameworku Protége. Výsledkem je logický systém obsahující důležité třídy, vlastnosti, vztahy a charakteristiky, který je schopný klasifikovat nové instance. K
\end{abstract}

\setlength{\parskip}{1ex plus 0.25ex minus 0.125ex}
  
\section{Úvod}
\label{sec:uvod}
Obor literatury je široký a komplikovaný, tato semestrální práce má za cíl jej dekomponovat do logických struktur a zpracovat ho do formy znalostní ontologie. Daná ontologie je nejprve modelována a následně zpracována v systému Progége. Výsledná ontologie má za cíl pokrývat vymezenou podmožinu  oboru literatury v souladu s omezeným záběrem této práce. 
V rámci vývoje ontologie byl použit systém \textit{Metoda 101}, který je iterativním procesem pro tvorbu znalostních ontologií.

\section{Metodika}
Pro vývoj ontologie byla vybrána \textit{Metoda 101} dle \cite{noy2001ontology}
, která je iterativním postupem pro vývoj ontologií . V rámci jednotlivých kroků byla vymezena doména, stanoveny hlavní důležité termíny, hierarchie tříd, vlastnosti, charakteristiky a nakonec byly vytvořeny instance a vztahy mezi nimi.
Takto vytvořený systém byl následně zadán do softwaru pro modelování ontologií Protége, který je výsledkem této práce.
Jmenné konvence jsou použity dle \cite{noy2001ontology} a jsou jednotné v rámci projektu. 

\section{Tvorba ontologie}
\subsection{Doména a rozsah záběru}
Démenou projektu jsou knihy a literatura. V rámci rozsahu byla tato doména omezena zejména na knihy a k nim vzahující se termíny. Doména knih byla dále omezena pouze na epiku, lyrická díla ani divadelní hry nebyly brány v úvahu. Kromě samotných knih do domény uvažujeme rovněž lidi, kteří s knihami běžně interagují. Tyto interakce jsou rovněž podrobeny zkoumání a relevantní interakce budou do projektu zahrnuty.

\subsection{Existující ontologie}
Vámci metodiky 101, byl vynechán krok zvážení existujících ontologií, vzhledem k rozsahu a zadání práce není možné žádnou ontologii převzít a cílem je vytvořit projekt nový. 

\subsection{Důležité termíny}
Mezi důležité termíny patří bezpochyby knihy. Hlavním termínem, který se ke knihám vztahuje byl určen člověk. Kromě člověka byly ke knihám stanoveny subsidiární termíny jazyk, autor, čtenář a formát. Tyto termíny se dále stanou vstupem pro definování tříd a jejich hierarchie, kde budou doplněny a dále strukturovány.

Celá ontologie by se tedy měla zaměřovat na knihy a člověka. Důležité je, v jaké formě se knihy k lidem dostávají, jak nejčastěji lidé interagují s knihami a jaké ostatní termíny dokáží vhodně popsat vztah knih a člověka.

Základní schopnost, zda člověk je schopen smysluplně 


\subsection{Třídy a jejich hierarchie}
Hierarchie tříd je základem každé znalostní ontologie. Pro výběr konkrétních tříd bylo přihlédnuto ke stanoveným důležitým termínům a ke zvolenému rozsahu domény. Mezi třídy byly stanoveny:
\begin{itemize}
    \item edice
    \item formát
    \item jazyk
    \item kniha (audiokniha, e-kniha, tištěná kniha)
    \item knihovna (databáze audioknih, e-shop, kamenná prodejna)
    \item období vydání
    \item osoba (autor, čtenář)
    \item série
    \item vydavatelství
    \item žánr
\end{itemize}
Jak je z třídní hierarchie patrné, pojmy autor a čtenář byly pro zachování korektnosti sloučeny pod entitu člověk a zmíněné důležité termíny byly doplněny o další vedlejší pojmy.
Jako jeden z hlavních prostředníků mezi člověkem a knihou je vybrána knihovna, kterou je myšleno médium, které je schopné poskytnout člověku knihu. Knihy se běžně vyskytují v různých formátech, tyto určují základní typ interakce mezi člověkem a danou knihou, jedná se tedy o jednu z vedlejších, avšak stále důležitých tříd. 

\subsection{Objektové vlastnosti}
V rámci objektových vlastnosti bylo vymezeno deset sloves, které formulují důležité vzahy mezi jednotlivými entitami v rámci domény. Mezi ně patří:
\begin{itemize}
    \item kniha \textbf{je dostupná} v knihovně
    \item kniha \textbf{je součást} knihovny
    \item kniha \textbf{je součástí} série
    \item kniha \textbf{je v} jazyce
    \item kniha \textbf{je vydána v} edici
    \item kniha \textbf{má} formát
    \item kniha \textbf{patří do} žánru
    \item autor \textbf{napsal} knihu
    \item vydavatelství \textbf{vydalo} knihu
    \item čtenář \textbf{čte} knihu
\end{itemize}

\subsection{Instance tříd}
Mezi instance byly vybírány konkrétní objekty, které splňují vlastnosti a vztahy jim přiřazených tříd. V rámci konstrukce třídní hierarchie byly některé původně zamýšlené třídy převedeny na instance, jedná se o:
\begin{itemize}
    \item formát - (audio, elektronický a tištěný formát)
    \item jazyk - (anglický a český)
    \item žánr - (fantasy, sci-fi a romantika)
\end{itemize}
Kromě zmíněných instací, které byly vytvořeny v rámci tvoření třídní hierarchie lze vytvářet instance nové, které model pomocí logických vztahů klasifikuje.

\subsection{Úplné definice}
Po vytvoření hierarchie tříd a objektových vlastností a instancí, které vznikly odtrhnutím z třídní hierarchie bylo možné přistoupit k tvorbě úplných definic některých tříd.
Mezi hlavní z úplných definic jsou řazeny:
\begin{itemize}
    \item tištěná kniha je ekvivalentní knize v nějakém jazyce kterou napsal nějaký autor a má tištěný formát
    \item audiokniha je ekvivalentní knize v nějakém jazyce kterou napsal nějaký autor a má audio formát
    \item databáze audioknih je ekvivalentní knihovně která má audio formát
    \item autor je ekvivalentní osobě která napsala nějakou knihu
    \item e-shop je je ekvivalentní knihovně která má elektronický formát

\end{itemize}
Po zadání úplných definic je systém schopen klasifikovat nové instance, které logickými vzahy spadají pod nějakou třídu. Např. po vytvoření instance \textit{Zaklínač 1: Poslední prání}, která je knihou, je v jazyce čeština a má tištěný formát jí Protége reasoner přiřadí do třídy tištěná kniha.
  

\section{Výsledky}
\label{sec:vysledky}
V rámci práce byl vytvořen systém znalostní ontologie, který obsahuje vymezenou doménu literatury, která se vzahuje zejména ke knihám. Systém obsahuje hierarchii tříd, vztahů, instancí a logických definic, které základním způsobem popisují tuto oblast lidského zaměření. Výstupem je zejména sestavený model, který je schopen klasifikovat nové instance na základě logických souvislostí a to především ohledně knih, jejich způsobů distribuce, autorů a čtenářů.

\section{Diskuze a závěr}
Získané výsledky ukazují, že systém je schopen klasifikovat základní instance z oboru literatury a knih. Kromě toho model přináší základní zpracování znalostní ontologie z tohoto oboru, na které je možné dále navázat v rámci případného rozšíření.

Práce využívá standardní přístup \textit{Metody 101} k modelování domény a zpracovává jej v moderním softwaru, který je volně dostupný.

Mezi hlavní limit práce patří úvodní omezení rozsahu domény , které se silně odráží v klasifikačních schopnostech modelu. Hierarchie tříd by se dala vzhledem k širokému záběru tohoto tématu volně rozšiřovat a celé toto odvětví by vhodněji popsal model výrazně větší.


\begin{thebibliography}{11}
\setlength{\parskip}{0ex}

	\bibitem{noy2001ontology}
Noy, N. F. and McGuinness, D. L. 
\textit{Ontology Development 101: A Guide to Creating Your First Ontology.} 
Technical Report, Stanford Knowledge Systems Laboratory, 2001.  
Available at: \url{https://protege.stanford.edu/publications/ontology_development/ontology101.pdf}.


\end{thebibliography}

\appendix
\section*{Přílohy}

\end{document}